\documentclass[11pt]{article}
\usepackage[margin=25mm]{geometry}
\usepackage{fontspec}
\setmainfont{Latin Modern Roman}
\usepackage{amsmath,amssymb,mathtools}
\usepackage{unicode-math}
\setmathfont{Latin Modern Math}
\usepackage{xurl,hyperref}
\hypersetup{colorlinks=true,urlcolor=blue}
\setcounter{secnumdepth}{0}
\setlength{\parindent}{0pt}
\setlength{\parskip}{5pt}
\setlength{\emergencystretch}{3em}
\newcommand{\sref}[2]{\hyperlink{source:#1:#2}{[S#2]}}
\newcommand{\cataloguescope}[1]{\paragraph{Recorded target.}#1}
\newcommand{\R}{\mathbb R}
\newcommand{\C}{\mathbb C}
\newcommand{\N}{\mathbb N}
\newcommand{\Z}{\mathbb Z}
\newcommand{\Q}{\mathbb Q}
\newcommand{\F}{\mathbb F}
\newcommand{\eps}{\varepsilon}
\begin{document}
% UnsolvedMath ID 30006352; source catalogue code OWR-14299511-001
\setcounter{section}{30006351}
\section{Reflection-Positive Two-Dimensional Quantum Field Theories with Boundary}\label{30006352}
{\small UnsolvedMath ID 30006352 · Mathematical Physics}
\subsection{Definitions and mathematical statement}

\paragraph{Shared notation.}$\mathbb N=\{1,2,\ldots\}$; the standard number fields have their usual meanings. All additional conventions are stated locally.


A von Neumann algebra is a weakly closed unital $*$-algebra of bounded operators. $O(2)$ describes rotations and reflections in two dimensions; a generalized action can involve higher or categorical data. Reflection positivity encodes positivity after a Euclidean reflection. The proposed flat ultraviolet setting has a vanishing energy--momentum tensor and product structures near the chosen boundary, together with reflection positivity and remote detectability. A bimodule carries commuting left and right algebra actions. The spaces $L^p(A)$ used for boundary angles need not be invertible bimodules.

\paragraph{Further primary-source questions.}Formulate the required noninvertible generalized $O(2)$ action. Determine what excludes concave boundary angles when the source's expression $L^{2\pi/\theta}(A)$ only makes sense for positive convex angles $\theta$ in its construction.

\paragraph{Catalogue wording retained for traceability.}
Identify the generalized $O(2)$-action on algebras whose fixed points are von Neumann algebras. Classify reflection-positive, remote-detectable, flat-space Lorentz-invariant two-dimensional quantum field theories with boundary: are von Neumann algebras the complete list, or do other examples exist?

\paragraph{Scope and source qualification.}The fixed-point claim and field-theory classification are both retained. The classification concerns the proposed flat ultraviolet and boundary axioms above; it is not a classification of arbitrary Lorentz-invariant two-dimensional quantum field theories.
\subsection{Short English statement}
Identify the generalized orthogonal symmetry and classify the specified two-dimensional boundary field theories, including possible examples beyond von Neumann algebras.
\subsection{Research attempt}
\paragraph{Model, date and investigated scope.}
OpenAI GPT-6 Astra (configured model \texttt{gpt-6-astra}), ultra reasoning; generated 11 October 2026. Verified a finite-matrix reflection pairing and a concrete noninvertible bimodule obstruction to ordinary group-action simplification.

\paragraph{Sources and literature check.}
Read the full supplied definition, catalogue wording, source qualifications and listed references for 30006352. No fresh status or novelty search was performed. The argument below proves its auxiliary statements directly and does not import an unverified external theorem.

\paragraph{A finite-algebra reflection test.} For \(A=M_n(\mathbb C)\), use the Hilbert space of matrices with \(\langle x,y\rangle=\operatorname{Tr}(x^*y)\). Left and right multiplication commute. The antilinear map \(Jx=x^*\) satisfies \(J^2=I\) and
\[
 J L_a J(x)=xa^*=R_{a^*}(x),\qquad
 \langle x,x\rangle=\operatorname{Tr}(x^*x)\geq0.
\]
Thus a reflection exchanges the two boundary actions and gives a positive pairing in this finite model. This elementary structure is consistent with the von Neumann-algebra examples but is not the full field-theory reflection-positivity axiom.
\paragraph{Why invertible actions are too restrictive.} Even for \(A=\mathbb C\), the bimodule \(H=\mathbb C^n\), \(n>1\), cannot have a finite-dimensional tensor inverse: \(\dim(H\otimes K)=n\dim K\ne1\) for every nonzero finite-dimensional \(K\). Thus replacing the source noninvertible generalized action by an ordinary group action through invertible bimodules discards basic candidates. Similarly the expression \(p=2\pi/\theta\) requires an actual positive angle in the stipulated construction; formally assigning negative \(p\) would fail on a nonzero matrix projection with zero eigenvalues and cannot define its usual Schatten norm.
\paragraph{First categorical/geometric gap.} These calculations do not define the generalized \(O(2)\) action, identify its fixed points, derive the angle composition law or exclude concave boundary geometries under the source conventions. They also do not prove remote detectability, a vanishing energy--momentum tensor or a classification of all allowed flat ultraviolet theories. The first gap is constructing the noninvertible action and its gluing/reflection coherence data; a positive trace pairing alone cannot supply them.

\paragraph{Verification and adversarial audit.}
Checked the antilinear reflection formula, trace positivity, tensor-dimension obstruction and the limited negative-exponent norm test; retained all boundary/classification qualifications. Model self-audit only; no independent human mathematical verification. Applied contribution prompt A before drafting and B afterwards. Checked the antilinear reflection formula, trace positivity, tensor-dimension obstruction and the limited negative-exponent norm test; retained all boundary/classification qualifications. First remaining gap: Construct the source noninvertible generalized O(2) action with angle, gluing and reflection coherence, before proving fixed points or field-theory completeness. The inherited source status is preserved; original-target completion remains zero. Additional GPT-6 Astra reviewer (ultra reasoning): Checked J²=I and J L\_a J(x)=x a* on the matrix Hilbert space, with positive Hilbert-Schmidt trace pairing. Checked that a finite-dimensional C-bimodule C\textasciicircum{}n with n>1 has no nonzero finite-dimensional tensor inverse because dimensions multiply. Checked the negative-Schatten-exponent warning using a projection with a nontrivial zero eigenspace; the draft does not equate this warning with a theorem excluding every concave geometry. Confirmed all restored source requirements remain explicit gaps: noninvertible generalized action, fixed points, angle/gluing coherence, remote detectability, vanishing stress tensor and the specified flat-ultraviolet classification.

\paragraph{Tools, process and contributions.}
Python 3 and local file tools used to read definitions and serialize manuscripts. Calculations are symbolic reasoning; no numerical experiment or theorem prover. Contribution prompt A was followed by prompt B. GPT-6 Astra with ultra reasoning drafted the research and performed its stated self-audit; any additional coordinator contribution is identified in the verification and audit. A GPT-6-based Codex coordinator assembled the manuscript and checked repository integration; its exact variant and reasoning configuration were not exposed. No human mathematical contribution, independent human proof check or formal proof certification is supplied. Complete derivations and the consulted source references are included here; private reasoning transcripts are not published.

\paragraph{Final status and first remaining gap.}
\textbf{UNRESOLVED}. Construct the source noninvertible generalized O(2) action with angle, gluing and reflection coherence, before proving fixed points or field-theory completeness.

\section{Completion Estimate}
\noindent\textbf{COMPLETION: 0\%}
This is an elementary auxiliary calculation or an explicit test of a restricted strategy. No new portion of the full source target is established.

\subsection{Sources}
\begin{itemize}
\item[\hypertarget{source:30006352:1}{S1}] ULAM.AI and the UnsolvedMath Contributors, UnsolvedMath, record 30006352 (OWR-14299511-001). Catalogue curation under CC BY 4.0; definitions and English exposition added locally.
\url{https://huggingface.co/datasets/ulamai/UnsolvedMath}.
\item[\hypertarget{source:30006352:2}{S2}] Subfactors and Applications (2025), Oberwolfach Reports, DOI 10.4171/owr/2025/34, pp. 1823--1886. Original report and mathematical context. Incorporated report material remains under the source CC BY-SA 4.0 terms unless otherwise noted.
\url{https://doi.org/10.4171/owr/2025/34}.
\end{itemize}
\label{end:30006352}
\end{document}
